% LTeX: language=de

\subsection{Grundlagen}

Grundlegende Aufgabenstellung: Bestimme die Anzahl (nicht die Wahrscheinlichkeit)

\begin{note}

Das \texthl{\textbf{allgemeine Zählprinzip}}: Anzahl der Ergebnisse eines mehrstufigen ZE erhält man durch multiplizieren der Anzahl der Ergebnisse jeder Stufe.\\[-1.0ex]

    \begin{Example}{Auswahl von Kleidungskombinationen}{}
        4 T-Shirts, 3 Hosen, 2 Paar Schuhe\\
        $\Rightarrow$ Wie viele verschiedene Kombinationen sind möglich?\\
        \raisebox{-\height}{
        \begin{forest}
        for tree={
            grow=south,                % grow down
            parent anchor=south, 
            child anchor=north,
            tier/.option=level, 
            l sep=8mm,             % level distance
            s sep=2mm,              % minimal sibling separation (can be auto-fitted)
            font=\footnotesize,
            rectangle, rounded corners, draw,
            align=center,
            edge+={-},
            edge label={node[midway, font=\footnotesize, fill=white, inner sep=1pt]{}} % default empty
        },
        [{}, draw=none  % root
        [T1, edge label={}, name=TA
            [\phantom{X}, draw=none, no edge
                [\phantom{X}, draw=none, no edge]
                [\phantom{X}, draw=none, no edge]
            ]
            [\phantom{X}, draw=none, no edge
                [\phantom{X}, draw=none, no edge]
                [\phantom{X}, draw=none, no edge]
            ]
        ]
        [T2, edge label={}, name=TB
            [H1, edge label={}, name=HA
                [S1, edge label={}, name=SA]
                [S2, edge label={}, name=SB]
            ]
            [H2, edge label={}, name=HB
                [\phantom{X}, draw=none, no edge]
                [\phantom{X}, draw=none, no edge]
            ]
            [H3, edge label={}, name=HC
                [\phantom{X}, draw=none, no edge]
                [\phantom{X}, draw=none, no edge]
            ]
        ]
        [T3, edge label={}, name=TC
            [\phantom{X}, draw=none, no edge
                [\phantom{X}, draw=none, no edge]
                [\phantom{X}, draw=none, no edge]
            ]
            [\phantom{X}, draw=none, no edge
                [\phantom{X}, draw=none, no edge]
                [\phantom{X}, draw=none, no edge]
            ]
        ]
        [T4, edge label={}, name=TD
            [\phantom{X}, draw=none, no edge
                [\phantom{X}, draw=none, no edge]
                [\phantom{X}, draw=none, no edge]
            ]
            [\phantom{X}, draw=none, no edge
                [\phantom{X}, draw=none, no edge]
                [\phantom{X}, draw=none, no edge]
            ]
        ]
        ]
        \node[anchor=west] at ($(TD.east)+(1.0,0)$) {
            \textbf{$4$}
        };
        \node[anchor=west] at ($(TD.east)+(1.0,-1.5)$) {
            \textbf{$3$}
        };
        \node[anchor=west] at ($(TD.east)+(1.0,-3.0)$) {
            \textbf{$2$}
        };
        \end{forest}
        }\\[2.0ex]
        $\Rightarrow$ $4 \cdot 3 \cdot 2 = 24$ verschiedene Kombinationen\\[1.0ex]
        Es ergeben sich sogenannte \texthl{\textbf{Tupel}}; \quad z. B. $\underbrace{\text{3-Tupel (T2 / H1 / S1)}}_{\footnotesize \text{\textcolor{myg}{Reihenfolge ist entscheidend; d. h. (T2 / T3 / S1)}}}$\\
    \end{Example}
\end{note}

\begin{Exercise}{S. 188/1}{}
\begin{enumerate}[leftmargin=*]
    \item [1.] {
        Serienbrief an 30 Personen; 16 Vorlagentexte; 3 Papiersorten\\[1.0ex]
        $N = 30 \cdot 16 \cdot 3 = 1440$
        }
    \item [2.]\begin{enumerate}[label=\alph*),leftmargin=*]
        \item[a)] {
            Nummernschild München mit 2 Buchstaben (inklusive Umlaute); 3-stellige Zahlen\\[1.0ex]
            $N = 29 \cdot 29 \cdot 10 \cdot 10 \cdot 10 = 841000$
            }
    \end{enumerate}
\end{enumerate}
\end{Exercise}
\pagebreak

\begin{note}
\begin{enumerate}[leftmargin=*]
    \item [1)] {
        \textcolor{myb}{\textbf{Ziehen \texthl{mit Zurücklegen} und mit \texthl{Beachtung der Reihenfolge}}}
        \begin{itemize}[leftmargin=*]
            \item alles was schon mal da war, darf nochmal verwendet werden
        \end{itemize}
        \begin{Example}{Zahlenschloss}{}
            Zahlenschloss mit 3 Rädern \'{a} 10 Ziffern (0-9)\\[1.0ex]
            \begin{tabularx}{0.5\linewidth}{X X}
                1. Rädchen: & 10 Möglichkeiten\\
                2. Rädchen: & 10 Möglichkeiten\\
                3. Rädchen: & 10 Möglichkeiten\\
            \end{tabularx}\\[1.5ex]
            $N = 10 \cdot 10 \cdot 10 = 10^3 = 1000$ verschiedene Kombinationen
        \end{Example}
        \begin{Corollary}{}{}
            k-maliges Ziehen aus n Möglichkeiten ergibt \textbf{$\highlight[yellow]{N = n^k}$} mögliche Ergebnisse
        \end{Corollary}
        }
\end{enumerate}
\end{note}

\begin{Exercise}{S. 194/3}{}
\begin{enumerate}[leftmargin=*]
    \item [3.]{
        \begin{itemize}[leftmargin=*]
            \item Passwort Mona: 5 Zeichen; Ziffern und Groß-/Kleinbuchstaben ohne Sonderzeichen
            \item Passwort Philip: 13 Zeichen; Kleinbuchstaben
        \end{itemize}
        \begin{enumerate}[label=\alph*),leftmargin=*]
            \item[a)] {
                \begin{itemize}[leftmargin=*]
                    \item Passwort Mona: \quad $N = (10 + 29 + 29)^5 = 68^5 = \SI{1453933568}{}$
                    \item Passwort Philip: \quad $N = 29^{13} = \SI{10260628710000000000}{}$
                \end{itemize}
                }
            \item[b)] {
                \begin{itemize}[leftmargin=*]
                    \item Passwort Mona: \quad $t = \dfrac{68^5}{\SI{1000}{\per\second} \cdot \SI{3600}{\second\per\hour}} = \SI{404}{\hour}$\\
                    \item Passwort Philip: \quad $t = \dfrac{29^{13}}{\SI{1000}{\per\second} \cdot \SI{3600}{\second\per\hour} \cdot \SI{24}{\hour\per\day} \cdot \SI{365}{\day\per a}} = \SI{7808697651}{a}$
                \end{itemize}
                }
        \end{enumerate}
        }
\end{enumerate}
\end{Exercise}

\begin{Exercise}{S. 194/6}{}
\begin{enumerate}[leftmargin=*]
    \item [6.]{
        \begin{itemize}[leftmargin=*]
            \item 1 Frage: \quad $N = 2^{6} = 64$
            \item 10 Fragen: \quad $N = 2^{6 \cdot 10} = 2^{60} = \SI{1.15e18}{}$
        \end{itemize}
    }
\end{enumerate}
\end{Exercise}
\pagebreak

\begin{note}
\begin{enumerate}[leftmargin=*]
    \item [2)] {
        \textcolor{myb}{\textbf{Ziehen \texthl{ohne Zurücklegen} und mit \texthl{Beachtung der Reihenfolge} (Variation)}}
        \begin{itemize}[leftmargin=*]
            \item Anzahl der Möglichkeiten verringert sich mit jedem Zug
        \end{itemize}
        \begin{Example}{Staffellauf}{}
            Beim Staffellauf werden 4 von 6 mögliche Läufer ausgewählt. Wie viele Aufstellungen sind möglich? \\[1.0ex]
            \begin{tabularx}{0.5\linewidth}{X X}
                1. Läufer: & 6 Möglichkeiten\\
                2. Läufer: & 5 Möglichkeiten\\
                3. Läufer: & 4 Möglichkeiten\\
                4. Läufer: & 3 Möglichkeiten\\
            \end{tabularx}\\[1.5ex]
            $N = 6 \cdot 5 \cdot 4 \cdot 3 = 360$ verschiedene Aufstellungen
        \end{Example}
        \begin{Definition}{Fakultät}{}
            Für eine natürliche Zahl $n$ ist das \texthl{\textbf{Produkt aller natürlichen Zahlen von 1 bis $n$}} die Fakultät von $n$ und wird mit $n!$ bezeichnet.\\[1.0ex]
            \[
                n! = n \cdot (n - 1) \cdot (n - 2) \cdots 3 \cdot 2 \cdot 1
            \]
            \[
                0! = 1
            \]
            \begin{Example}{Beispiel Fakultät}{}
                $5! = 5 \cdot 4 \cdot 3 \cdot 2 \cdot 1 = 120$           
            \end{Example}
        \end{Definition}
        \begin{align*}
            N &= n \cdot (n - 1) \cdot (n - 2) \cdots (n - k  + 1) = \dfrac{n!}{(n - k)!} \\[1.0ex]
              &= \dfrac{n \cdot (n - 1) \cdot (n - 2) \cdots (n - k + 1) \cdot (n - k) \cdots}{(n - k) \cdot (n - k - 1) \cdots}
        \end{align*}
        \begin{Corollary}{}{}
            k-maliges Ziehen aus n-elementiger Menge ergibt 
            \begin{equation*}
                \highlight[yellow]{N = \dfrac{n!}{(n - k)!}} 
            \end{equation*}
            mögliche Ergebnisse.\\[1.0ex]
            Sonderfall \texthl{\textbf{Permutation}}: \textbf{Alle} n Elemente werden verteilt (d. h. $k = n$); \quad $\highlight[yellow]{N = n!}$
        \end{Corollary}
        }
\end{enumerate}
\end{note}

\begin{Exercise}{S. 194/2}{}
\begin{enumerate}[leftmargin=*]
    \item [2.]{
        4 Personen verteilt auf 4 Parkplätze\\[1.0ex]
        $N = 4! = 4 \cdot 3 \cdot 2 \cdot 1 = 24$
    }
\end{enumerate}
\end{Exercise}

\begin{note}
\begin{enumerate}[leftmargin=*]
    \item [3)] {
        \textcolor{myb}{\textbf{Ziehen \texthl{ohne Zurücklegen} und ohne \texthl{Beachtung der Reihenfolge} (Kombination)}}
        \begin{itemize}[leftmargin=*]
            \item Reihenfolge ist nicht entscheidend
            \item Anzahl der Möglichkeiten verringert sich mit jedem Zug
        \end{itemize}
        \begin{Example}{Lottospiel}{}
            Beim Lottospiel werden 3 von 7 mögliche Zahlen ausgewählt. Wie viele Kombinationen sind möglich? \\[1.0ex]
            Falls die Reihenfolge eine Rolle spielen würde (Variation):\\
            \begin{equation*}
                N = \dfrac{7!}{(7 - 3)!} = 7 \cdot 6 \cdot 5 = 210 \text{ verschiedene Tipps}
            \end{equation*}
            Davon sind ohne Beachtung der Reihenfolge der Tipps gleich:\\
            \begin{equation*}
                \text{z. B.}: \quad (1 / 2 / 3) = (1 / 3 / 2) = (2 / 1 / 3) = (2 / 3 / 1) = (3 / 1 / 2) = (3 / 2 / 1)
            \end{equation*}
            Wie viele Möglichkeiten gibt es, um 3 Zahlen auf 3 Plätze zu verteilen?\\[1.0ex]
            \begin{equation*}
                N = 3! = 3 \cdot 2 \cdot 1 = 6 \text{ Möglichkeiten}
            \end{equation*}
            $\Rightarrow$ Von den 210 möglichen Tipps (falls die Reihenfolge relevant wäre) sind \underline{jeweils 6 Tipps identisch}, falls die Reihenfolge keine Rolle spielt.\\[1.0ex]
            \begin{equation*}
                N = \dfrac{7!}{(7 - 3)! \cdot 3!} = \dfrac{7 \cdot 6 \cdot 5}{3 \cdot 2 \cdot 1} = 35 \text{ verschiedene Tipps}
            \end{equation*}
        \end{Example}
            \begin{Corollary}{}{}
            k-maliges Ziehen aus n-elementiger Menge ergibt 
            \begin{equation*}
                \underbrace{\highlight[yellow]{N = \dfrac{n!}{(n - k)! \cdot k!} = \binom{n}{k}}}_{\textcolor{myg}{\text{Binomialkoeffizient; } k \text{ aus } n}}
            \end{equation*}
            mögliche Ergebnisse.\\[1.0ex]
            Der Taschenrechner hat eine spezielle Taste für den Binomialkoeffizienten: \quad \texttt{nCr}
        \end{Corollary}
        \begin{Example}{Beispiel Binomialkoeffizient}{}
            $\displaystyle \binom{8}{3} = \dfrac{8!}{(8 - 3)! \cdot 3!} = \dfrac{8 \cdot 7 \cdot 6}{3 \cdot 2 \cdot 1} = 56$
        \end{Example}
        }
\end{enumerate}
\end{note}

\pagebreak

\begin{Exercise}{S. 194/4}{}
\begin{enumerate}[leftmargin=*]
    \item [4.]{
        4 Kinder; 2 Karten\\[1.0ex]
        $\displaystyle \binom{4}{2} = \dfrac{4!}{(4 - 2)! \cdot 2!} = 6$
    }
\end{enumerate}
\end{Exercise}

\begin{Exercise}{S. 194/5}{}
\begin{enumerate}[leftmargin=*]
    \item [5.]{
        15 Schüler; 4 Karten; n = 15; k = 4\\[1.0ex]
        \begin{enumerate}[label=\alph*),leftmargin=*]
            \item[a)] {
                $\displaystyle \binom{15}{4} = \dfrac{15!}{(15 - 4)! \cdot 4!} = 1365$\\
                }
            \item[b)] {
                $\dfrac{15!}{(15-4)!} = 32760$\\
                }
            \item[c)] {
                $15^{4} = 50625$
                }
        \end{enumerate}
        }
\end{enumerate}
\end{Exercise}

\begin{Exercise}{S. 195/11}{}
\begin{enumerate}[leftmargin=*]
    \item [11.]{
        \begin{itemize}
            \item Belgien: \quad $\displaystyle \binom{42}{6} = \SI{5245786}{}$
            \item Sweden: \quad $\displaystyle \binom{35}{7} = \SI{6724520}{}$
        \end{itemize}
        $\Rightarrow$ Gewinnchance in Belgien sind besser\\[1.0ex]
        Weitere Kriterien:
        \begin{itemize}
            \item Gewinnsumme
            \item Teilgewinnsummen bei z. B. nur 4 richtigen
            \item Kosten pro Tipp
        \end{itemize}
    }
\end{enumerate}
\end{Exercise}

\pagebreak

\begin{note}
\begin{enumerate}[leftmargin=*]
    \item [4)] {
        \textcolor{myb}{\textbf{Mississippi-Problem}}\\[1.0ex]
        Bestimmen Sie, wie viele verschiedene Anordnungen man aus den Buchstaben des Wortes MISSISSIPPI bilden kann (alle sollen verwendet werden).\\[1.0ex]
        Annahme: alle 11 Buchstaben sind unterscheidbar $\Rightarrow 11!$ mögliche Anordnungen\\[1.0ex]
        Verschiedene Möglichkeiten aber nur falls alle Buchstaben unterscheidbar sind.\\
        Es gibt:
        \begin{itemize}[leftmargin=*]
            \item $4!$ Möglichkeiten I zu tauschen
            \item $4!$ Möglichkeiten S zu tauschen
            \item $2!$ Möglichkeiten P zu tauschen
        \end{itemize}
        Sind gleiche Buchstaben nicht zu unterscheiden, so gibt es:
        \begin{equation*}
            \dfrac{{11!}}{4! \cdot 4! \cdot 2!} = \dfrac{11 \cdot 10 \cdot 9 \cdot 8 \cdot 7 \cdot 6 \cdot 5}{4 \cdot 3 \cdot 2 \cdot 4 \cdot 3 \cdot 2 \cdot 2} = 34650
        \end{equation*}
        }
\end{enumerate}
\end{note}

\begin{Exercise}{S. 195/13}{}
\begin{enumerate}[leftmargin=*]
    \item [13.]{
        \begin{enumerate}[label=\alph*),leftmargin=*]
            \item[a)] {
                HOF\\[0.5ex]
                $N = \dfrac{3!}{1! \cdot 1! \cdot 1!} = 6$\\
                }
            \item[b)] {
                PASSAU\\[0.5ex]
                $N = \dfrac{6!}{1! \cdot 2! \cdot 2! \cdot 1!} = 180$\\
                }
            \item[c)] {
                REGENSBURG\\[0.5ex]
                $N = \dfrac{10!}{2! \cdot 2! \cdot 2!} = 453600$\\
                }
        \end{enumerate}
        }
\end{enumerate}
\end{Exercise}